%%%%%%%%%%%%
%
% $Autor: Wings $
% $Datum: 2019-03-05 08:03:15Z $
% $Pfad: TemplateSensor $
% $Version: 4250 $
% !TeX spellcheck = en_GB/de_DE
% !TeX encoding = utf8
% !TeX root = filename 
% !TeX TXS-program:bibliography = txs:///biber
%
%%%%%%%%%%%%

% Structure
\chapter{A4988-Schrittmotortreiber}

\section{Einleitung}

Für die Ansteuerung der Schrittmotoren im Projekt wird der Schrittmotortreiber A4988 verwendet. 
Der A4988 ist ein integrierter DMOS-Microstepping-Treiber für bipolare Schrittmotoren und wird in Positioniersystemen, CNC-Systemen und ähnlichen Anwendungen eingesetzt \cite{Allegro:A4988}.

Im Projekt übernimmt der A4988 die Leistungsansteuerung der Achsmotoren. 
Dadurch können die Bewegungen des Heizdrahts definiert und wiederholbar ausgeführt werden.

\section{Grundlagen}

Schrittmotoren gehören zu den elektrischen Antrieben, bei denen die Drehbewegung in einzelne Schritte unterteilt wird. 
Dadurch können Positionen gezielt angefahren werden, ohne dass zwingend ein zusätzliches Positionsmesssystem erforderlich ist \cite{Isermann:2008}.

Für die Ansteuerung eines bipolaren Schrittmotors muss die Stromrichtung in den Motorwicklungen umgeschaltet werden. 
Diese Aufgabe wird durch leistungselektronische Schaltungen übernommen \cite{Zach:2022}. 
Der A4988 enthält hierfür integrierte H-Brücken und eine interne Stromregelung.

Die Ansteuerung erfolgt über die Signale STEP und DIR. 
Mit jeder steigenden Flanke am STEP-Eingang bewegt der Treiber den Motor um einen Schritt beziehungsweise Mikroschritt weiter. 
Das DIR-Signal bestimmt die Drehrichtung \cite{Allegro:A4988}.

\section{Spezifischer Treiber}

Im Projekt wird ein A4988-kompatibles Treibermodul verwendet. 
Der eigentliche Treiber-IC stammt von Allegro MicroSystems. 
Laut Herstellerdatenblatt unterstützt der A4988 Voll-, Halb-, Viertel-, Achtel- und Sechzehntelschrittbetrieb \cite{Allegro:A4988}. 
Durch Microstepping kann die Bewegung der Achsen gleichmäßiger ausgeführt werden.

Der A4988 besitzt eine integrierte PWM-Stromregelung. 
Dadurch wird der Strom durch die Motorwicklungen begrenzt. 
Dies ist wichtig, da der Motor mit einer höheren Versorgungsspannung betrieben werden kann, ohne die Wicklungen zu überlasten.

Zusätzlich verfügt der Treiber über Schutzfunktionen gegen Überstrom, Kurzschluss, Unterspannung und Übertemperatur \cite{Allegro:A4988}.

\subsection{Installation}

Der A4988 wird direkt auf das CNC Shield V3 gesteckt. 
Bei Verwendung des CNC Shields sind die Anschlüsse für STEP, DIR und ENABLE bereits den jeweiligen Achsen zugeordnet.

Vor dem Anschluss des Motors muss die Strombegrenzung eingestellt werden. 
Außerdem müssen die Masse des Mikrocontrollers und die Masse der externen Motorspannungsversorgung miteinander verbunden werden.

\section{Spezifikation}

Die technischen Daten des verwendeten A4988 wurden dem Herstellerdatenblatt entnommen \cite{Allegro:A4988}.

\begin{itemize}
	\item Motorspannung: 8\,V bis 35\,V
	\item Logikspannung: 3\,V bis 5{,}5\,V
	\item maximaler Ausgangsstrom: bis $\pm$2\,A
	\item Schrittmodi: Vollschritt, Halbschritt, Viertelschritt, Achtelschritt und Sechzehntelschritt
	\item Steuerung über STEP- und DIR-Eingang
	\item integrierte PWM-Stromregelung
	\item Schutzfunktionen gegen Übertemperatur, Unterspannung und Überstrom
\end{itemize}

Der maximale Motorstrom ergibt sich aus:

\[
I_{TripMAX} = \frac{V_{REF}}{8 \cdot R_S}
\]

Dabei beschreibt \(R_S\) den Sense-Widerstand und \(V_{REF}\) die Referenzspannung \cite{Allegro:A4988}.

\section{Bibliothek}

\subsection{Beschreibung}

Für den A4988 wird keine zusätzliche Arduino-Bibliothek benötigt. 
Die Ansteuerung erfolgt direkt über digitale Ausgänge des Mikrocontrollers. 
Im späteren Betrieb des CNC-Hot-Wire-Foam-Cutters erzeugt die GRBL-Firmware die notwendigen STEP- und DIR-Signale.



\subsection{Funktionen}

Zur direkten Steuerung des A4988 werden Standardfunktionen der Arduino-IDE verwendet. 
Mit \texttt{pinMode()} werden die verwendeten Pins als Ausgänge konfiguriert. 
Mit \texttt{digitalWrite()} werden STEP-, DIR- und ENABLE-Signale ausgegeben. 
Die Funktion \texttt{delayMicroseconds()} bestimmt die Zeit zwischen den Schrittimpulsen und damit die Drehgeschwindigkeit.

\section{Beispiel}

In diesem Abschnitt wird ein einfaches Beispiel beschrieben, das die grundlegende Funktion des A4988 zeigt. 
Der Mikrocontroller erzeugt Schrittimpulse, die vom Treiber in Motorbewegungen umgesetzt werden.

\subsection{Versuchsaufbau}

Der A4988 wird mit einem bipolaren Schrittmotor verbunden. 
Die Motorspannung wird über ein externes Netzteil bereitgestellt. 
Der Arduino erzeugt die STEP- und DIR-Signale. 
Bei Verwendung des CNC Shield V3 wird der Treiber auf den Steckplatz der gewünschten Achse gesteckt.

\subsection{Funktionsweise des Beispiels}

Im Beispiel wird zunächst die Drehrichtung festgelegt. 
Anschließend werden am STEP-Eingang periodische Impulse ausgegeben. 
Jeder Impuls bewegt den Motor um einen Schritt oder Mikroschritt weiter.

\subsection{Code}

Für die Dokumentation werden drei Programme verwendet. 
Das erste Programm prüft die Grundfunktion des Treibers. 
Das zweite Programm zeigt eine einfache projektbezogene Anwendung. 
Das dritte Programm dient als vollständiger Funktionstest.

\subsection{Dateien}

\begin{itemize}
	\item A4988SimpleStep.ino
	\item A4988FoamCutterAxis.ino
	\item A4988DriverTest.ino
\end{itemize}

\section{Kalibrierung}

Vor dem Betrieb muss die Strombegrenzung des A4988 eingestellt werden. 
Die Einstellung erfolgt über ein Potentiometer auf dem Treibermodul.

Eine zu hohe Strombegrenzung führt zu starker Erwärmung und kann eine thermische Abschaltung verursachen. 
Eine zu niedrige Strombegrenzung kann Schrittverluste verursachen. 
Für den CNC-Hot-Wire-Foam-Cutter wird die Strombegrenzung so eingestellt, dass der Motor zuverlässig läuft und der Treiber nicht übermäßig heiß wird.

\section{Einfacher Code}

Der einfache Code dient zur Überprüfung der Grundfunktion des A4988. 
Der Motor wird kontinuierlich in eine Richtung bewegt. 
Dadurch können Verdrahtung, Spannungsversorgung, STEP-Signal und Treiberfunktion überprüft werden.

\begin{code}[H]
	\captionof{code}{Einfacher Funktionstest des A4988-Schrittmotortreibers}
	\label{Code:A4988SimpleStep}
	\ArduinoExternal{}{../../Code/A4988SimpleStep/A4988SimpleStep.ino}
\end{code}

\section{Einfache Anwendung}

Die einfache Anwendung simuliert eine Achsbewegung des CNC-Hot-Wire-Foam-Cutters. 
Die Achse bewegt sich langsam in Schneidrichtung und anschließend schneller zurück. 
Dadurch wird der praktische Einsatz des A4988 im Projekt demonstriert.

\begin{code}[H]
	\captionof{code}{Einfache Achsbewegung des CNC-Hot-Wire-Foam-Cutters}
	\label{Code:A4988FoamCutterAxis}
	\ArduinoExternal{}{../../Code/A4988FoamCutterAxis/A4988FoamCutterAxis.ino}
\end{code}

\section{Tests}

\subsection{Einfacher Funktionstest}

Beim einfachen Funktionstest wird überprüft, ob der Motor grundsätzlich auf STEP-Signale reagiert. 
Dieser Test erfolgt mit dem Programm aus Abschnitt \ref{Code:A4988SimpleStep}. 
Dabei werden die Verdrahtung, die Versorgungsspannung und die Drehrichtung kontrolliert.

\subsection{Test aller Funktionen}

Beim vollständigen Test werden Drehrichtung, Geschwindigkeit, ENABLE-Signal und Temperaturverhalten überprüft. 
Der Motor wird in beide Richtungen bewegt. 
Zusätzlich wird der Treiber kurz deaktiviert und anschließend wieder aktiviert. 
Während des Tests wird beobachtet, ob der Motor Schritte verliert oder ob sich der Treiber übermäßig erwärmt.

\begin{code}[H]
	\captionof{code}{Testprogramm für den A4988-Schrittmotortreiber}
	\label{Code:A4988DriverTest}
	\ArduinoExternal{}{../../Code/A4988DriverTest/A4988DriverTest.ino}
\end{code}

\section{Weiterführende Literatur}

\begin{itemize}
	\item Isermann, Rolf: \textit{Mechatronische Systeme}. Springer, 2008. \cite{Isermann:2008}
	\item Zach, Franz: \textit{Leistungselektronik}. Springer Vieweg, 2022. \cite{Zach:2022}
	\item Allegro MicroSystems: \textit{A4988 DMOS Microstepping Driver with Translator and Overcurrent Protection}. Datenblatt, 2022. \cite{Allegro:A4988}
\end{itemize}